% Abstract.
% All figures are from the rebuild and rerun of 2026-10-20; the evidence
% bundle is archived in reports.zip and summarized in Appendix B.

This project implements a mini search engine in C++17 for \emph{The Complete
Works of William Shakespeare}. It covers the four required tasks---word
frequency statistics and stop-word identification, a stemmed inverted index,
word and phrase queries (with a Boolean extension), and an experiment on query
thresholds---and addresses the bonus question of 500,000 files and 400,000,000
distinct words with measured results and a clearly marked extrapolation.

The corpus is extracted from 1,211 HTML files mirrored from the MIT Shakespeare
site and reduced to 195 documents (37 plays, 154 sonnets, and 4 long poems),
with 152,632 lines and 930,911 words; the default tokenizer produces 866,626
body tokens. Three IETF RFCs form a second corpus that tests whether the engine
is independent of the Shakespeare format.

For stop-word selection, a two-feature tf/df rule narrows the candidate band,
and a length-normalized rate entropy separates words that are spread evenly
across documents from words that are merely frequent. The resulting
three-feature rule derives 39 stop words, the same count as the manually
curated list, with an overlap of 35 words. The rule also adapts to the RFC
corpus: it derives a different 58-word list there and correctly excludes the
RFC 2119 keywords \code{must}, \code{may}, and \code{should} because their
distributions are not flat enough.

For the index, Porter stemming maps 25,705 word forms to 14,947 stems. The
formal index contains 14,937 terms, 117,704 postings, and 556,866 positions.
The first persistence format (v1, uncompressed) occupies 3,685,655 bytes; the
second format (v2, delta plus variable-length coding) reduces this to
1,725,694 bytes, or 46.82\% of v1. Splay-tree, hash-table, and \code{std::map}
dictionaries produce byte-identical indexes. A disk-based B+ tree sidecar
(237,568 bytes) reduces cold-query I/O to 0.225\% of a full load, and a
SPIMI-style block build finishes in 17 runs and produces an index identical to
the in-memory build.

The query program supports word, phrase, and Boolean queries and exposes four
thresholds: \code{max-df}, \code{min-df}, \code{min-match}, and \code{top-x}.
On a matrix of 15 queries and 201 measurements, the thresholds change query
semantics rather than merely precision: removing a query term can turn a
three-term conjunction into a single-term query and raise the result count from
30 to 90, while phrase queries ignore the term-side thresholds.

Correctness is supported by 18 CTest tests (8 Catch2 executables, 46 test
cases, and 85,155 assertions) and by independent Python implementations that
reproduce the word-count CSV, rebuild the full index, parse both persistence
formats, and check 303 phrase queries. The bonus experiments use a synthetic
corpus of 500,000 documents and 20,000,000 tokens: block construction needs
1.01\,GiB of peak memory versus 2.00\,GiB for the in-memory build, and both
produce byte-identical indexes. The B+ tree dictionary is measured from
100,000 to 50,000,000 keys at a stable 15.100 bytes per key, with height 4 at
50,000,000 keys; extrapolating the measured fanout to 400,000,000 words gives
about 6.04\,GB on disk, height 4, and four page reads (about 16\,KiB) per
lookup.

\vspace{0.8em}
\noindent\textbf{Keywords:} inverted index; word frequency; stop words; rate
entropy; Porter stemming; splay tree; B+ tree; query thresholds;
reproducibility
